\documentclass[12pt,a4paper]{article}
\usepackage[UTF8]{ctex}
\usepackage{amsmath, amssymb}
\usepackage{geometry}
\usepackage{booktabs}
\usepackage{listings}
\usepackage{xcolor}
\usepackage{tcolorbox}
\usepackage{enumitem}
\usepackage{hyperref}
\usepackage{float}
\usepackage{tikz}
\usetikzlibrary{arrows.meta, positioning}

\geometry{left=2.5cm, right=2.5cm, top=2.5cm, bottom=2.5cm}

\lstset{
    language=Python,
    basicstyle=\ttfamily\small,
    keywordstyle=\color{blue}\bfseries,
    commentstyle=\color{gray},
    stringstyle=\color{red},
    showstringspaces=false,
    breaklines=true,
    frame=single,
    numbers=left,
    numberstyle=\tiny\color{gray},
    tabsize=4,
}

\title{\textbf{Q3 代码架构与运行逻辑完全指南}\\[6pt]
\large VLSI 布图规划设计 --- 最小可行死区比例的二分搜索}
\author{2026 年第七届"华数杯"大学生数学建模竞赛 B 题}
\date{2026/08}

\begin{document}
\maketitle
\tableofcontents
\newpage

% ============================================================
\section{项目总体架构}
% ============================================================

\subsection{目录结构}

\begin{lstlisting}[language=bash]
The_7th_Huashu_Cup.../
|-- cpp_solver/                  # C++ 共享核心（Q2 模式 + --feas-only 判定）
|   `-- src/
|       |-- floor_plan.hpp       # B*-Tree 拓扑 + Skyline 解码 + 解析器
|       |-- sa.hpp               # 两阶段退火（run 找可行 + run2 精修）
|       `-- main.cpp             # CLI（q2 模式 / --feas-only / --seed）
|-- common/
|   `-- verify.py                # 独立合法性复核
|-- q3/                          # 问题 3 客制层（<-- 本文件）
|   |-- q3_bisect.py             # 二分搜索 + 并行多种子判定 + 双向确认
|   |-- q3_solver.py             # 驱动：二分 + d* 处完整求解 + 汇总出图
|   |-- output/
|   |   |-- q3_metrics.csv       # d*/side/HPWL/confirm 记录/合法性
|   |   |-- q3_<chip>_bisect.csv # 二分轨迹（可审计）
|   |   |-- q3_<chip>.rpt/log    # d* 处最优布局 + 收敛轨迹
|   `-- visualization/
|       |-- plot_bisect.py       # 二分轨迹图（区间收缩/可行标记/d* 标注）
|       |-- plot_{floorplan,convergence,snapshots}.py
|       `-- figs/<时间戳>/<chip>/  # 每芯片 12 张图（含 bisect 轨迹）
`-- paper/
\end{lstlisting}

\subsection{执行顺序}

\begin{tcolorbox}[colback=blue!5, colframe=blue!60, title=执行顺序]
\begin{enumerate}
    \item \texttt{cd cpp\_solver \&\& make}（编译 bin/main）
    \item \texttt{conda activate math}
    \item \texttt{python q3/q3\_solver.py --repeats 8}
          （二分 + d* 处多起点求解 + 汇总 + 出图）
\end{enumerate}
\end{tcolorbox}

\subsection{核心设计理念}

\textbf{问题本质}：随死区比例 $d$ 减小，轮廓 $W=H=\lceil\sqrt{\Sigma A(1+d)}\rceil$
同步缩小，存在可行布局的概率单调下降——最小可行 $d^*$ 是"可行性随 $d$ 单调"
的二分搜索问题。

\textbf{判定可靠性}：SA 是启发式，"找不到可行解"可能是搜索不足（假阴性）而非
真不可行。为此：① 每比例\textbf{并行多种子判定}（任一可行即可行）；
② 收敛后\textbf{双向确认}（$d^*$ 本身必须可行，$d^*-\varepsilon$ 处多种子
必须全不可行）——保证 $d^*$ 是"验证过的最小稳定点"。

% ============================================================
\section{Q3 求解流程}
% ============================================================

\begin{figure}[H]
\centering
\begin{tikzpicture}[
  node distance=1.0cm,
  box/.style={draw, rounded corners=2pt, fill=blue!8, align=center,
              text width=7.2cm, minimum height=0.9cm},
  dec/.style={draw, rounded corners=2pt, fill=orange!15, align=center,
              text width=7.2cm, minimum height=0.9cm},
  term/.style={draw, rounded corners=8pt, fill=green!12, align=center,
               text width=7.2cm, minimum height=0.9cm},
  arrow/.style={-{Stealth[length=3mm]}, thick},
]
\node[term] (start) {二分区间 $[0, 0.15]$（0.15 已知可行）};
\node[box, below=of start] (mid) {$mid=(lo+hi)/2$，理论下界断言 $side^2\ge\Sigma A$};
\node[box, below=of mid] (check) {并行多种子判定（粗区间 2、精区间 3 种子）\\ \texttt{bin/main q2 --feas-only --seed s}};
\node[dec, below=of check] (branch) {任一判定可行？};
\node[box, below left=1.0cm and -0.5cm of branch] (hi) {$hi=mid$};
\node[box, below right=1.0cm and -0.5cm of branch] (lo) {$lo=mid$};
\node[dec, below=of hi] (term) {$hi-lo<10^{-4}$？};
\node[box, below=of term] (confirm) {双向确认：$d^*$ 可行、$d^*-\varepsilon$ 多种子全不可行\\（否则上移/下移重验，最多 12 步）};
\node[box, below=of confirm] (full) {在 $d^*$ 处完整求解（多起点 8 轮取 HPWL 最优）};
\node[term, below=of full] (out) {产出：\texttt{q3\_metrics.csv} + 每芯片 12 图};
\draw[arrow] (start) -- (mid);
\draw[arrow] (mid) -- (check);
\draw[arrow] (check) -- (branch);
\draw[arrow] (branch) -- node[left] {是} (hi);
\draw[arrow] (branch) -- node[right] {否} (lo);
\draw[arrow] (lo) -- ++(0,-0.6) -| (check);
\draw[arrow] (hi) -- (term);
\draw[arrow] (term) -- node[left] {否} (mid);
\draw[arrow] (term) -- node[right] {是} (confirm);
\draw[arrow] (confirm) -- (full);
\draw[arrow] (full) -- (out);
\end{tikzpicture}
\caption{Q3 二分搜索 + 双向确认流程}
\end{figure}

% ============================================================
\section{二分搜索算法（q3\_bisect.py）}
% ============================================================

\subsection{判定调用（--feas-only）}

共享核心提供可行性判定模式：
\begin{lstlisting}[language=bash]
./bin/main q2 0.5 <blocks> <nets> <pl> <rpt> <dead> --feas-only --seed <s>
# 输出 3 行：feasible(0/1) / bbox W H / hpwl
\end{lstlisting}

\begin{itemize}
    \item 只跑 \texttt{run()} 可行性阶段，找到首个可行解即返回（早停）。
    \item reset 上限保证单次判定有界（$N/16+1$ 次温度重置）。
    \item 固定种子 $\Rightarrow$ 判定可复现。
\end{itemize}

\subsection{并行多种子判定}

单种子判定存在假阴性（差种子搜索不足）。每比例并行 $K$ 个独立种子
（粗区间 $K=2$、精区间 $K=3$），\textbf{任一可行即可行}：

\begin{equation}
\text{feasible}(d) = \bigvee_{s=1}^{K} \text{feasible}_s(d),
\qquad
P(\text{假阴性}) \approx \prod_{s=1}^{K} P(\text{假阴性}_s)
\end{equation}

判定种子确定性与互异：\texttt{seed = BASE + chip\_idx*10000 + tag*100 + i}。

\subsection{二分与终止}

\begin{itemize}
    \item 区间 $[0, 0.15]$（0.15 即 Q2 基准，已知可行）。
    \item 终止：$hi - lo < 10^{-4}$（约 11 次迭代，精度 $1.4\times10^{-7}$ 内）。
    \item 每步记录轨迹（iter/lo/hi/mid/side/bbox/feasible/耗时）到
          \texttt{q3\_<chip>\_bisect.csv}（可审计）。
    \item 理论下界全程断言：$side^2 \ge \Sigma A$ 恒成立。
\end{itemize}

\subsection{双向确认（正确性关键）}

二分收敛的 $d_{bisect}$ 可能受判定波动影响，做\textbf{双向确认}：

\begin{enumerate}
    \item 在 $d$ 处做 $K_{verify}=3$ 种子判定，\textbf{必须可行}；
          若不可行 $\Rightarrow$ 上移 $d \mathrel{+}= \varepsilon$ 重验。
    \item 在 $d-\varepsilon$ 处做 3 种子判定，\textbf{必须全不可行}；
          若有可行 $\Rightarrow$ 下移 $d \mathrel{-}= \varepsilon$ 重验。
    \item 循环直至两点同时满足（最多 12 步），记录每步结果。
\end{enumerate}

最终 $d^*$ 满足：\texttt{d\_feas=True} 且 \texttt{below\_infeas=True}
（$d^*-\varepsilon$ 处 $[False,False,False]$）——即
\textbf{验证过的最小稳定点}，写入 \texttt{confirm\_steps} 可审计。

% ============================================================
\section{Phase：d* 处完整求解（q3\_solver.py）}
% ============================================================

\begin{itemize}
    \item 在 $d^*$ 处调用完整 Q2 求解（run + run2 $\times 2$），
          多起点 8 轮并行取 HPWL 最优（复用 q2 逻辑）。
    \item 只接受可行轮次（bbox $\le$ side）；保留 best 的 rpt/log。
    \item 独立复核：无重叠 / 尺寸保持 / 不越界（common/verify.py）。
    \item 产出：\texttt{q3\_metrics.csv}（d*/side/HPWL/confirm/合法性）
          + 每芯片 12 张图（9 快照 + 收敛 + 布图 + bisect 轨迹）。
\end{itemize}

% ============================================================
\section{判定可靠性讨论}
% ============================================================

\subsection{为什么需要并行多种子判定？}

早期版本（单种子判定）二分结果受种子运气影响：同一比例不同种子判定结果
可相互矛盾（0.14209 判不可行、0.14206 判可行），导致 $d^*$ 被推高
（n300 达 0.142）。并行多种子判定将假阴性概率按种子数指数下降，
n300 $d^*$ 从 0.142 改善到 0.125。

\subsection{为什么需要双向确认？}

二分只保证"最后一个判可行的 mid 为上界"，不保证 $d^*-\varepsilon$ 处
真不可行。双向确认把"可行/不可行"两个方向都验证到稳定，
$d^*$ 的语义是"在当前搜索强度下验证过的最小可行比例"——
论文表述严谨（不声称超越搜索能力的真最优）。

\subsection{与理论下界的关系}

$d$ 的理论下界为 0（$side^2 \ge \Sigma A$ 恒满足）；实际最小 $d$ 由
打包效率决定。搜索能力内的 $d^*$ 是保守估计（可能略高于真值），
但确认记录保证其"验证过"性质。

% ============================================================
\section{测试与验证}
% ============================================================

\begin{itemize}
    \item \textbf{理论下界断言}：每次判定 $side^2 \ge \Sigma A$。
    \item \textbf{confirm 记录}：\texttt{d\_feas} 与 \texttt{below\_infeas}
          双条件满足才算通过（\texttt{d\_confirmed\_minimum=True}）。
    \item \textbf{轨迹审计}：\texttt{q3\_<chip>\_bisect.csv} 每步完整记录。
    \item \textbf{最终布局复核}：Python 独立验证（无重叠/尺寸/越界）。
    \item C++ 测试：feas-only 正反例、run 快照、HPWL oracle（1135 断言）。
\end{itemize}

% ============================================================
\section{实际运行结果}
% ============================================================

\subsection{最小死区比例（并行判定 + 双向确认）}

\begin{table}[H]
\centering
\begin{tabular}{lccccccc}
\toprule
芯片 & $d^*$ & side & HPWL & 迭代 & 确认步数 & 确认通过 & 合法性 \\
\midrule
n100 & 0.117307 & 448 & 293671 & 11 & 2 & 是 & 合法 \\
n200 & 0.108977 & 442 & 544534 & 11 & 4 & 是 & 合法 \\
n300 & 0.124885 & 555 & 826685 & 11 & 6 & 是 & 合法 \\
\bottomrule
\end{tabular}
\end{table}

\begin{itemize}
    \item 全部芯片 \texttt{d\_confirmed\_minimum=True}：
          最终 $d^*-\varepsilon$ 处 3 种子判定全不可行
          （\texttt{below\_results=[False,False,False]}）。
    \item n300 确认过程含 4 次上移（$d^*$ 处判定波动），最终在
          0.124885 稳定——确认记录完整可审计。
    \item $d^*$ 处 HPWL 大于 Q2（轮廓更紧、搜索更困难），合法且合理。
\end{itemize}

% ============================================================
\section{文件依赖关系}
% ============================================================

\begin{lstlisting}
q3_solver.py
  |-- q3_bisect.py               (二分 + 判定 + 确认)
  |-- cpp_solver/bin/main        (q2 模式 + --feas-only + --seed)
  |-- data/raw/{n100,n200,n300}.{blocks,nets,pl}
  |-- common/verify.py
  `-- q3/visualization/plot_{bisect,floorplan,convergence,snapshots}.py
\end{lstlisting}

% ============================================================
\section{运行指南}
% ============================================================

\subsection{环境要求}

\begin{itemize}
    \item g++（\texttt{-std=c++17}）；conda 环境 \texttt{math}。
    \item TeX Live 2026（xelatex + ctex，编译本指南）。
\end{itemize}

\subsection{完整运行流程}

\begin{lstlisting}[language=bash]
cd cpp_solver && make
cd ..
conda activate math
python q3/q3_solver.py --repeats 8   # 二分 + d* 求解 + 汇总 + 出图
\end{lstlisting}

\subsection{常见问题}

\begin{enumerate}
    \item \textbf{d* 偏保守}：判定强度与搜索深度相关——可提高判定种子数
          （\texttt{PRECISE\_SEEDS}）或判定搜索强度（参数优化阶段）。
    \item \textbf{确认失败}：\texttt{d\_confirmed\_minimum=False} 时
          结果不可交付——增大确认步数或判定种子后重跑。
    \item \textbf{复现}：所有判定/求解使用固定种子公式
          （\texttt{BASE + chip\_idx*10000 + tag*100 + i}），可精确复现。
\end{enumerate}

% ============================================================
\section{合规清单}
% ============================================================

\begin{itemize}
    \item 原始数据只读；二分/判定/求解全部种子显式管理（可复现）。
    \item \texttt{d\_confirmed\_minimum} 为交付硬性前提（三芯片均为 True）。
    \item 判定波动与保守性如实记录（第 5 节），未作任何美化。
    \item 论文表述建议："在当前搜索强度下验证过的最小可行死区比例"。
\end{itemize}

\end{document}
